\section{Intuitionistic Synthesis of Ranking Functions}

This section proposes domain-stable sequences and
constructs a ranking function from a given SC graphs 
under SCT condition and fan-in-free condition
by using intuitionistic logic.

\subsection{Domain-Stable Sequences}

This section proposes domain-stable sequences.
%
We assume a SC call graph in this section.

In this paper, we call an SC graph a {\em relation}.
When $A \subset X \times Y$,
we call $A$ an {\em ordinary relation} on $X$ and $Y$.
Then a relation is represented by $(f,(S,P),g)$ where
$f,g$ are function symbols and $S,P$ are ordinary relations.

\begin{Def}[Injective]\rm
An ordinary relation $R$ is defined to be {\em injective}
if $x R z$ and $y R z$ imply $x=y$.
\end{Def}

A relation $(f,(S,P),g)$
corresponds to a fan-in free SC graph \cite{Benamram09}
when $S \cup P$ is injective.

\begin{Def}[Composition of SC graphs]\rm
The composition $R_1\ldots R_k$ of a given SC path $(R_1,\ldots,R_k)$ is
defined as
%
$(f_1,(S_1,P_1)\ldots(S_k,P_k),f_{k+1})$
%
where
%
$R_i=(f_i,(S_i,P_i),f_{i+1})$ for all $i \in [1,k]$,
%
and we define
%
$(S_1,P_1)\ldots(S_k,P_k) = (S,P)$
%
where
%
$S = S_1\ldots S_k$ (ordinary composition of ordinary relations), and
%
$P = (S_1 \cup P_1)\ldots (S_k \cup P_k) - S$.

We write $i(R_1\ldots R_i)j$ for $i (S \cup P) j$.
%
We call a composition of a SC path a {\em SC composition}.
%
We will often write $R$ for a SC composition.
%
A sequence $(R_1,\ldots,R_k)$ of SC compositions 
is defined to be {\em composable} if
$R_i = (f_i,-,f_{i+1})$ for all $i \in [1,k-1]$.
%
Note that a SC path is a special case of a composable sequence of SC compositions.
%
A composition of a composable sequence of SC compositions
is defined in a similar way to a composition of a SC path.

A SC composition $R$ is called {\em injective} if
the relation $S \cup P$ is injective
where
the composition $R = (-,(S,P),-)$.
A sequence $(R_1,\ldots,R_k)$ of SC compositions is defined to be {\em injective} if
each $R_i$ is injective for $i \in [1,k]$.
As a special case,
a SC path is defined to be {\em injective} in a similar way to 
a sequence of SC compositions.
A SC graph is defined to be {\em injective} if all edges in the graph are
injective.

For a SC composition $R$, we define
%
$\dom(R)$ as $\{x \ |\ R = (-,(S,P),-), x(S \cup P)y \}$ and
%
$\codom(R)$ as $\{y \ |\ R = (-,(S,P),-), x(S \cup P)y \}$.
%
For a finite set $X$, we write $|X|$ for the cardinal of $X$.
%
Note that $\dom(R)$ and $\codom(R)$ for a SC path $R$ is 
defined as a special case of this definition.
\end{Def}

\begin{Def}[Domain-Stable]\rm
A sequence $(R_1,\ldots,R_k)$ of SC compositions
is defined to be {\em domain-stable}
if 
the sequence is composable, and
for all $i \in [1,k-1]$, 
we have $\dom(R_i) = \dom(R_iR_{i+1})$.
\end{Def}

\begin{Lemma}[Cardinal of Domains]\label{lemma:cardinal}
If $|\dom(R)|=|\dom(RS)|$, then $(R,S)$ is domain-stable.
\end{Lemma}

{\em Proof.}
By definition of $\dom$, we have $\dom(R) \supseteq \dom(RS)$.
Since $|\dom(R)|=|\dom(RS)|$, we have $\dom(R) = \dom(RS)$.
$\Box$

\begin{Def}\rm
For a given composable sequence $R_1,\ldots,R_n$ of SC compositions,
a sequence $R'_1,\ldots,R'_l$ of SC compositions
is defined to be a {\em partial composition} of $R_1,\ldots,R_n$
with head indexes $a_1,\ldots,a_l$
if 

- there is some sequence $R_{a_1},\ldots,R_{b_1},R_{a_2},\ldots,R_{b_2},\ldots, R_{a_l},\ldots,R_{b_l}$
which is the same sequence as $R_1,\ldots,R_n$
where $a_i \le b_i$ for all $i \in [1,l]$,

- $R_i' = R_{a_i}\ldots R_{b_i}$ for all $i \in [1,l]$.
\end{Def}

\begin{Def}\rm
We define a {\em value assignment} to be a map from $N$ to $N$.
Note that we will use a value assignment $\alpha$ so that
$\alpha(i)=m$ means that the $i$-th variable has a value $m$.

We write a sequence $(v_1,\ldots,v_n)$ of natural numbers for
a value assignment $\alpha$ such that
$\alpha(i)=v_i$ for all $i \in [1,n]$ and
$\alpha(i)=0$ otherwise.

For a SC composition $R$ and value assignments $\alpha$ and $\alpha'$,
we define $\alpha R \alpha'$ to hold if for all $k,l \in N$,
$\alpha(k) \ge \alpha'(l)$ for $k S l$, and
$\alpha(k) > \alpha'(l)$ for $k P l$,
where
$R_i = (f,(S,P),g)$.
\end{Def}

Note that $\exists \alpha_2(\alpha_1 R_1 \alpha_2 R_2 \alpha_3)$ 
implies $\alpha_1 R_1R_2 \alpha_3$, but
$\alpha_1 R_1R_2 \alpha_3$ does not imply 
 $\exists \alpha_2(\alpha_1 R_1 \alpha_2 R_2 \alpha_3)$ 
since 
$\alpha_1(1) =2 > \alpha_2(1) = \alpha_3(1) =1$ and
$1 P_1 1 P_2 1$ may happen.

\begin{Def}[Value Assignment Along Sequence and Domain-Sum Value]\rm
We define $\sigma$ to be a {\em value assignment sequence along} a composable sequence $\pi$ of
SC compositions 
if $\pi$ is $R_1,\ldots, R_n$,
and $\sigma$ is $\Vec v_1,\ldots,\Vec v_n$,
and 
%
 $\Vec v_i R_i \Vec v_{i+1}$ for all $i \in [1,n-1]$.

For a SC composition $R$ and a value assignment $\alpha$,
we define a {\em domain-sum value} $\Val(R,\alpha) \in N$
as $\Sigma \{ \alpha(i) \ |\ i \in \dom(R) \}$.
\end{Def}

\begin{Lemma}\label{lemma:A}
if $R_1,R_2,R_3$ is composable,
and $\Vec v_1,\Vec v_2$ is along $R_1,R_2$, 
and $\Vec v_2,\Vec v_3$ is along $R_2,R_3$, 
then
$\Vec v_1,\Vec v_3$ is along $R_1R_2,R_3$.
\end{Lemma}

{\em Proof.}
Since $\exists j(i R_1 j R_2 k)$ implies $i R_1R_2 k$,
we have $\Vec v_1 (R_1R_2) \Vec v_3$.
$\Box$

\begin{Prop}[Transform to Domain-Stable Sequence]\label{prop:domain-stable}
(1) For a composable sequence $R_1,\ldots,R_k,R$ of SC compositions,
if $R_1,\ldots,R_k$ is domain-stable
and $R_k,R$ is not domain-stable, then
the following sequcne $R_1,\ldots,R_{i-1},R'$ is
domain-stable and a partial composition of $R_1,\ldots,R_k,R$
with head redexes $1,\ldots,i-1,i$ and $R_{i},R_{i+1}\ldots R_k R$ is
not domain-stable.

- Take $i$ to be the smallest $i \in [1,k]$ among $i$'s such that
$R_i, R_{i+1} \ldots R_k R$ is not domain-stable.
It exists since the sequence $R_k,R$ is not domain-stable.
Let
\[
R' = R_i\ldots R_k R.
\]

(2) For any composable sequence $R_1,\ldots,R_k$ of SC compositions,
the following sequence $R_1',\ldots,R_l'$ 
is domain-stable and
a partial composition of 
$R_1,\ldots,R_k$ with head indexes $d_1,\ldots,d_l$.
Moreover,
if $\Vec v_1,\ldots,\Vec v_k$ is along $R_1,\ldots,R_k$,
then $\Vec v'_1,\ldots,\Vec v'_l$ is along $R'_1,\ldots,R'_l$,
where $\Vec v'_i$ is defined as $\Vec v_{d_i}$.

For each $0 < n \le k$, we define a sequence $R^{(n)}$ of SC compositions and 
a sequence $H^{(n)}$ of numbers by induction on $n$.
First we define
\[
R^{(1)} = (R_1), \ \ \
H^{(1)} = (1).
\]
We define $R^{(n+1)}$ and $H^{(n+1)}$ from $R^{(n)}$ and $H^{(n)}$ as follows.
Let
\[
R^{(n)} = (R_1',\ldots,R_m'), \ \ \
H^{(n)} = (n_1,\ldots,n_m).
\]
Case 1. If $R_1',\ldots,R_m',R_{n+1}$ is domain-stable,
then we define
\[
R^{(n+1)} = (R_1',\ldots,R_m',R_{n+1}), \ \ \
H^{(n+1)} = (n_1,\ldots,n_m,n+1).
\]
Case 2. If $R_1',\ldots,R_m',R_{n+1}$ is not domain-stable.

By (1), there is some $i \in [1,m]$ such that
$R'_1,\ldots,R'_{i-1}, R_i'\ldots R_m'R_{n+1}$ is domain-stable
and $R_i',R_{i+1}'\ldots R_m'R_{n+1}$ is not domain-stable.
Define
\[
R^{(n+1)} = (R_1',\ldots,R_{i-1}',R_i'R'_{i+1}\ldots R'_m R_{n+1}), \ \ \
H^{(n+1)} = (n_1,\ldots,n_{i}).
\]

We define $R_1',\ldots,R_l'$ as $R^{(k)}$
and define $n_1,\ldots,n_l$ as $H^{(k)}$.
\end{Prop}

{\em Proof.}
(1) It immediately follows from the definition of $i$.

(2) We can show $R^{(n)}$ is domain-stable for all $n \in [1,k]$
by induction on $n$.
If $\Vec v_1,\ldots,\Vec v_k$ is along $R_1,\ldots,R_k$,
then by induction on $l$,
$\Vec v'_1,\ldots,\Vec v'_l$ is shown to be along $R'_1,\ldots,R'_l$.
$\Box$

The construction of a domain-stable sequence 
in this proposition can be explained as follows.
For a given $R_1,\ldots,R_k$,
we handle them from the left to the right.
Assume we have already constructed a domain-stable sequence $R_1',\ldots,R_m'$
from the original sequence $R_1,\ldots,R_n$ for $n<k$.
Next we handle $R_{n+1}$. If $R_1',\ldots,R_m',R_{n+1}$ is
domain-stable, we take the resulting sequence to be
$R_1',\ldots,R_m',R_{n+1}$.
Otherwise
we handle $R_1',\ldots,R_m'$ from the right to the left.
If $R_1',\ldots,R_{m-1}', R_m'  R_{n+1}$ is domain-stable,
we stop.
Otherwise, if $R_1',\ldots,R_{m-1}'  R_m'  R_{n+1}$ is domain-stable,
we stop. Continue it until
$R_1',\ldots,R'_{i-1},R'_i \ldots R_{m-1}' R_m'  R_{n+1}$ is domain-stable.
Then we take it to be the resulting domain-stable sequence
for a sequence $R_1,\ldots,R_n$.
After handling the whole sequence $R_1,\ldots,R_k$,
the resulting sequence $R_1',\ldots,R_l'$ is 
some maximal partial composition of a given sequence $R_1,\ldots,R_k$
such that $R_1',\ldots,R_l'$ is domain-stable.
Here a head index is taken from the index of the head element of each group.

\begin{Lemma}\label{lemma:domain-stable-increase}
If
a composable sequence $R_1,\ldots,R_k$ of SC compositions
is domain-stable and injective, then
$|\dom(R_1)| \le \ldots \le |\dom(R_k)|$.
\end{Lemma}

{\em Proof.}
Let $R = R_i$ and $S = R_{i+1}$.

Since $R,S$ is domain-stable, we have $\dom(RS) = \dom(R)$.
Hence for each $x \in \dom(R)$ there are $y,z$ such that $x R y S z$.
Define a function $f: \dom(R) \to \dom(S)$ such that $f(x)=y$.
Since $f$ is injective, we have $|\dom(R)| \le |\dom(S)|$.
$\Box$

\begin{Lemma}[Grouping by Cardinal of Domains]\label{lemma:cardinal-constant}
For any 
domain-stable injective
sequence $R_1,\ldots,R_k$ of SC compositions,
there are $a_1,b_1,\ldots, a_l,b_l \in [1,k]$ such that
$R_{a_1},\ldots,R_{b_1},R_{a_2},\ldots,R_{b_2},\ldots,R_{a_l},\ldots,R_{b_l}$
is the same sequence as $R_1,\ldots,R_k$ such that

- $|\dom(R_i)| = |\dom(R_j)|$ for all $m \in [1,l]$ and
all $a_m \le i \le j \le b_m$,

- $|\dom(R_{a_i})| < |\dom(R_{a_{i+1}})|$ for all $i \in [1,l-1]$.

\end{Lemma}

{\em Proof.}
By Lemma \ref{lemma:domain-stable-increase}.
$\Box$

\begin{Lemma}\label{lemma:bijection}
If $R,S$ is a domain-stable sequence
and $|\dom(S)|=|\dom(R)|$
and $R$ is injective,
then $R:\dom(R) \to \dom(S)$ is a bijection.
\end{Lemma}

Its proof is given in the appendix.

\subsection{Synthesis of Ranking Functions}

This section constructs a ranking function from a given SC graphs.

\begin{Lemma}\label{lemma:maxlen}
If a SC call graph satisfies SCT condition,
then
for a domain-stable injective sequence $R_1,\ldots, R_n$ of SC compositions
and a value assignment sequence $\Vec v_1,\ldots,\Vec v_n$ along $R_1,\ldots, R_n$
such that $|\dom(R_i)|=k$ for all $i \in [1,n]$ and
$\Val(R_i,\Vec v_i)$ is the same for all $i \in [1,n]$,
we have
$n \le |V|{K \choose k}$
where $V$ is the set of vertexes of the SC call graph and
$K$ is the maximum arity in all function names.
\end{Lemma}

{\em Proof.}
By Lemma \ref{lemma:bijection},
$R_i:\dom(R_i) \to \dom(R_{i+1})$ is a bijection for all $i \in [1,n-1]$.

Let $R_i=(a_i,(S_i,P_i),a_{i+1})$ for all $i \in [1,n]$.

Since $\Val(R_i,\Vec v_i)$ is the same for all $i \in [1,n]$,
$P_i = \emptyset$ for all $i \in [1,n-1]$.

For any $i <j \in [1,n]$, we have $\dom(R_i) \ne \dom(R_j)$ or $a_i \ne a_j$.
We can show it as follows.
Assume $\dom(R_i) = \dom(R_j)$ and $a_i = a_j$
for some $i < j \in [1,n]$ to show contradiction.
We consider an infinite composable sequence $(R_i,\ldots,R_{j-1})^\omega$.
By applying SCT condition to it,
there is a progressing trace of some tail of it.
Hence
there are some $l \in [i,j-1]$ such that $R_l$ 
is progressing, which contradicts
with $P_l = \emptyset$.

Hence $(a_1,\dom(R_1)),\ldots, (a_n,\dom(R_n))$ are different.
The number of $a_i$ is $\le |V|$.
The number of $\dom(R_i)$ is $\le {K \choose k}$.
Hence $n \le |V|{K \choose k}$.
$\Box$

\begin{Def}[Weight]\label{def:weight}\rm
For a SC call graph that satisfies SCT condition,
a composable sequence $\pi$ of injective SC compositions, and 
a value assignment sequence $\sigma$ along $\pi$,
we define the weight $W(\pi, \sigma) \in ((N\cup\{\infty\}) \times N)^{K}$ 
where $K$ is the maximum arity for all function names
as follows.

Let $\pi$ be $R_1,\ldots, R_n$ and
$\sigma$ be $\Vec v_1,\ldots,\Vec v_n$.

By Proposition \ref{prop:domain-stable} (2), we define
a domain-stable sequence $R_1',\ldots,R_l'$ 
of partial SC compositions of $R_1,\ldots, R_n$ with 
head indexes $d_1,\ldots,d_l$.

We define $\Vec v'_i$ as $\Vec v_{d_i}$ for all $i \in [1,l]$.
Then by Proposition \ref{prop:domain-stable} (2), 
the value assignment sequence $\Vec v'_1,\ldots,\Vec v'_l$ is
along $R_1',\ldots,R_l'$.

By Lemma \ref{lemma:cardinal-constant},
we have
$R'_{a_1},\ldots,R'_{b_1},R'_{a_2},\ldots,R'_{b_2},\ldots, 
R'_{a_r},\ldots,R'_{b_r}$ such that

- $R'_{a_1},\ldots,R'_{b_1},R'_{a_2},\ldots,R'_{b_2},\ldots,
R'_{a_r},\ldots,R'_{b_r}$
is the same sequence as $R'_1,\ldots,R'_l$,

- $|\dom(R'_i)| = |\dom(R'_j)|$ for all $k \in [1,r]$ and
all $a_k \le i \le j \le b_k$,

- $|\dom(R'_{a_i})| < |\dom(R'_{a_{i+1}})|$ for all $i \in [1,r-1]$.

Fix $k \in [1,K]$. We define $W_k$.

Case 1. There is no $i \in [1,r]$ such that $|\dom(R'_i)| = k$.

We define $W_k$ as $(\infty,0)$.

Case 2. There is some $i \in [1,r]$ such that $|\dom(R'_i)| = k$.

Take $i \in [1,r]$ such that $|\dom(R'_{a_i})| = k$.

We define $W_k$ as $(\Val(R'_{b_i}, \Vec {v'}_{b_i}),|V|{K \choose k}-s)$
where $V$ is the set of vertexes of the SC call graph,
$K$ is the maximum arity for all function names,
and $s$ is the maximum $s$ such that
$b_i-s+1 \in [a_i,b_i]$ and
$\Val(R'_{b_i}, \Vec {v'}_{b_i}) = \Val(R'_{b_i-s+1}, \Vec {v'}_{b_i-s+1})$.
Note that by Lemma \ref{lemma:maxlen}, we have $|V|{K \choose k} \ge s$.

Note also that $s$ is the maximum length of the sequence from $R'_{b_i}$
to the left
such that each element has the same value by $\Val$.
\end{Def}

\begin{Def}\rm
We write $>_\Lex$ for the lexicographic order on $((N\cup\{\infty\}) \times N)^K$
for some $K$.
\end{Def}

\begin{Lemma}\label{lemma:stable-extension}
For a SC call graph that satisfies SCT condition,
domain-stable injective SC sequences $(\pi,\pi')$ and $(\pi,R)$ of SC compositions and 
a value assignment sequence $(\sigma,\sigma')$ along $(\pi,\pi')$ and
a value assignment sequence $(\sigma,\Vec v)$ along $(\pi,R)$,
if $|\pi'|\ne 0$ and $\pi'=(R',-)$ and $|\dom(R)|<|\dom(R')|$ or $|\pi'|=0$,
then
\[
W((\pi,\pi'),(\sigma,\sigma'))
>_\Lex
W((\pi,R),(\sigma,\Vec v)).
\]
\end{Lemma}

{\em Proof.}
Let $\pi=(R_1,\ldots,R_n)$ and $\sigma = (\Vec v_1,\ldots,\Vec v_n)$.
Let $k=|\dom(R)|$.
Let $v = \Val(R,\Vec v)$.
Let $W((\pi,\pi'),(\sigma,\sigma')) = W = (W_1,\ldots,W_K)$ and
$W((\pi,R),(\sigma,\Vec v)) = W' = (W'_1,\ldots,W'_k)$.

We have $W_j = W'_j$ for $j \in [1,k-1]$, since $|\pi'|\ne 0$ and $|\dom(R)|<|\dom(R')|$
or $|\pi'|=0$.

By Lemma \ref{lemma:domain-stable-increase}, 
we have $|\dom(R_{n})| \le k$.

Let $L$ be $|V|{K \choose k}$ 
where $V$ is the set of vertexes of the SC call graph
and $K$ is the maximum arity for all function names.

Case 1. $k=|\dom(R_{n})|$.

By Lemma \ref{lemma:bijection},
$
R_n:\dom(R_n) \to \dom(R)
$
is a bijection.
Hence we have $\Val(R_n,\Vec {v}_{n}) \ge v$.

Case 1.1. $\Val(R_n,\Vec {v}_{n}) = v$.

Let $W_k = (v,b)$.
Then $W'_k = (v, b-1)$.
Hence $W_k >_\Lex W'_k$.
Since $W_j = W'_j$ for $j \in [1,k-1]$,
we have $W >_\Lex W'$.

Case 1.2. $\Val(R_{n},\Vec {v}_{n}) > v$.

Let $W_k = (\Val(R_{n},\Vec {v}_{n}),b)$.
Then $W'_k = (v, L)$.
Hence $W_k >_\Lex W'_k$.
Since $W_j = W'_j$ for $j \in [1,k-1]$,
we have $W >_\Lex W'$.

Case 2. $k>|\dom(R_{n})|$.

Then $W_k = (\infty,0)$ since 
$|\pi'|\ne 0$ and $\pi'=(R',-)$ and $|\dom(R)|<|\dom(R')|$ or $|\pi'|=0$.
We have $W'_k = (v, L)$.
Hence $W_k >_\Lex W'_k$.
Since $W_j = W'_j$ for $j \in [1,k-1]$,
we have $W >_\Lex W'$.
$\Box$

\begin{Prop}\label{prop:decrease-pi}
For a SC call graph that satisfies SCT condition, and
a composable SC sequence $(\pi,R)$ of SC compositions and 
a value assignment sequence $(\sigma,\Vec v)$ along $(\pi,R)$,
we have $W(\pi,\sigma) >_\Lex W((\pi,R),(\sigma,\Vec v))$.
\end{Prop}

{\em Proof.}
Let
$\pi' = (\pi,R)$ and
$\sigma' = (\sigma,\Vec v)$.
Let $W(\pi,\sigma) = (W_1,\ldots,W_K)$ and
$W(\pi',\sigma') = (W'_1,\ldots,W'_K)$.

By Proposition \ref{prop:domain-stable} (2) for $\pi$, we define
a domain-stable sequence $R_1',\ldots,R_l'$ 
of partial SC compositions of $\pi$ with 
head indexes $d_1,\ldots,d_l$.
%
Define $\Vec v'_i$ as $\Vec v_{d_i}$ for all $i \in [1,l]$.
%
Then by Proposition \ref{prop:domain-stable}, 
the value assignment sequence $\Vec v'_1,\ldots,\Vec v'_l$ is
along $R'_1,\ldots,R'_l$.

Case 1. $R'_1,\ldots,R'_l,R$ is not domain-stable.

By Proposition \ref{prop:domain-stable} (1),
there is some $i$ such that $R'_1,\ldots,R'_{i-1},R'_i\ldots R'_lR$
is domain-stable and $R'_i,R'_{i+1}\ldots R'_lR$ is not domain-stable.
Let
\[
R' = R'_i\ldots R'_lR.
\]
Since $R'_i,R'_{i+1}\ldots R'_lR$ is not domain-stable and Lemma \ref{lemma:cardinal},
we have $|\dom(R')|<|\dom(R'_i)|$.

By applying Lemma \ref{lemma:stable-extension}
to $R'_1,\ldots,R'_l$ with $\Vec v'_1,\ldots,\Vec v'_l$ and
$R'_1,\ldots,R'_{i-1},R'$ with 
$\Vec v'_1,\ldots,\Vec v'_{i-1},\Vec v'_i$ with
$|\dom(R')|<|\dom(R'_i)|$,
we have
\[
W((R'_1,\ldots,R'_l),(\Vec v'_1,\ldots,\Vec v'_l)) >_\Lex
W((R'_1,\ldots,R'_{i-1},R'),
(\Vec v'_1,\ldots,\Vec v'_{i-1},\Vec v'_i)).
\]
Since
$
W(\pi,\sigma) = W((R'_1,\ldots,R'_l),(\Vec v'_1,\ldots,\Vec v'_l))
$ and
$
W(\pi',\sigma') = W((R'_1,\ldots,R'_{i-1},R'),
(\Vec v'_1,\ldots,\Vec v'_{i-1},\Vec v'_i))
$
by definition, we have 
$W(\pi,\sigma) >_\Lex W(\pi',\sigma')$.

Case 2. $R'_1,\ldots,R'_l,R$ is domain-stable.

By applying Lemma \ref{lemma:stable-extension}
to $R'_1,\ldots,R'_{l}$ with $\Vec v'_1,\ldots,\Vec v'_l$ and
$R'_1,\ldots,R'_{l},R$ with $\Vec v'_1,\ldots,\Vec v'_l,\Vec v$,
we have
\[
W((R'_1,\ldots,R'_{l}), (\Vec v'_1,\ldots,\Vec v'_l)) >_\Lex
W((R'_1,\ldots,R'_{l},R), (\Vec v'_1,\ldots,\Vec v'_l,\Vec v)).
\]
Since
$
W(\pi,\sigma) = W((R'_1,\ldots,R'_{l}), (\Vec v'_1,\ldots,\Vec v'_l))
$
and
$
W(\pi',\sigma') = 
W((R'_1,\ldots,R'_{l},R), (\Vec v'_1,\ldots,\Vec v'_l,\Vec v))
$
by definition, we have 
$W(\pi,\sigma) >_\Lex W(\pi',\sigma')$.
$\Box$

\begin{Lemma}\label{lemma:shorten}
For a domain-stable sequence $\pi$ of SC compositions
and a value assignment sequence $\sigma$ along $\pi$,
if $|\pi| > C_0$, then
we have a partial composition $\pi'$ of $\pi$ with head indexes 
$a_1,\ldots,a_l$ such that
$\sigma'$ is along $\pi'$ 
and $|\pi'| \le C_0$
and 
\[
W(\pi,\sigma)=W(\pi',\sigma')
\]
where
$C_0 = |V|(2^K-1)+1$ and
$(\sigma')_{i-1} = (\sigma)_{a_i-1}$ for all $i \in [1,l]$.
\end{Lemma}

Its proof is given in the appendix.

\begin{Def}[Value Assignment Sequence Minimally Along SC Compositions]\rm
We define $\Vec v$ to be the {\em minimum value assignment 
for an injective SC composition $R$ and $\Vec v'$} if the following hold.
\[
-\ \Vec v(i) = \max\{ \Vec v'(j)+1,\Vec v'(k) \ |\ i P j, i S k \} \hbox{\ if for some $j$, \ } i (S\cup P) j, \\
-\ \Vec v(i) = 0 \hbox{\ otherwise},
\]
where $R=(-,(S,P),-)$.
We write $\Vec v \min(R) \Vec v'$ when
$\Vec v$ to be the {\em minimum value assignment 
for $R$ and $\Vec v'$}.

$\sigma$ is defined to be {\em minimally along} $\pi$ if
for all $i \in [1,n-1]$, $\Vec v_i$ is the minimum value assignment
for $R_i$ and $\Vec v_{i+1}$,
where $\pi=(R_1,\ldots,R_n)$ and $\sigma=(\Vec v_1,\ldots,\Vec v_n)$.
\end{Def}

Note that for given a composable sequence $\pi$ of SC compositions 
and $\Vec v$, 
a value assignment sequence $\sigma$ minimally along $\pi$ and $\Vec v$ uniquely exists.

\begin{Th}[Intuitionistic Synthesis of Ranking Function]\label{th:rankingfunc}
For an injective SC call graph that satisfies SCT condition,
we define
\[
\rho(f,\Vec x) = \min\{ W(\pi,\sigma) \in ((N\cup\{\infty\}) \times N)^{K} \ |\ 
\\ \qquad
\Last(\pi) = (-,-,f),
\Last(\sigma) \min(\Last(\pi)) \Vec x,
\pi \hbox{\ a domain-stable SC compositions},
\\ \qquad
\sigma \hbox{\ a value assignment sequence minimally along $\pi$},
|\pi| \le C_0
\}
\]
where 
$K$ is the maximum arity in the SC call graph, and
$\Last(x)$ returns the last element of the sequence $x$, and
$V$ is the set of vertexes of the SC call graph,
and $C_0 = |V|(2^K-1)+1$.

If $s \to^G s'$ for some SC graph $G$,
then $\rho(s) >_\Lex \rho(s')$.
\end{Th}

{\em Proof.}
Assume $(f,\Vec x) \to^G (g,\Vec y)$ to show
$\rho(f,\Vec x) >_\Lex \rho(g,\Vec y)$.

Let $R$ be $(f,-,g)$.

Define $\pi$ and $\sigma$ as those that give the minimum
in the body of the definition of $\rho(f,\Vec x)$.
Then
$
\rho(f,\Vec x)=W(\pi,\sigma)
$
and
$
\Last(\pi) = (-,-,f)
$
and
$
\Last(\sigma) \min(\Last(\pi)) \Vec x
$
and
$\pi$ is a domain-stable SC compositions and
$\sigma$ is a value assignment sequence minimally along $\pi$ and
$|\pi| \le C_0$.

Define $\pi'$ and $\sigma'$ by
$
\pi'=(\pi,R)
$
and
$
\sigma'=(\sigma,\Vec x).
$
Then we have
$
\Last(\pi') = (-,-,g)
$
and
$
\sigma' \hbox{\ along\ } \pi'
$
and
$\Last(\sigma') \Last(\pi') \Vec y.$

By Proposition \ref{prop:decrease-pi},
we have
$W(\pi,\sigma) >_\Lex W(\pi',\sigma')$.

By definition of weights,
$
W(\pi',\sigma')=W(\pi'',\sigma'')
$
and
$\pi''$ is domain-stable
and $\Last(\pi'')=(-,-,g)$
and $\Last(\sigma'') \Last(\pi'') \Vec y$
and $\sigma''$ is along $\pi''$,
where $\pi''$ and the head redexes $d_1,\ldots,d_l$ are 
defined to be obtained from $\pi'$ by
Proposition \ref{prop:domain-stable} (2) and
we define $(\sigma'')_{i-1} = (\sigma')_{d_i-1}$ for all $i \in [1,l]$.

By Lemma \ref{lemma:shorten},
we have a partial composition $\pi'''$ of $\pi''$ and $\sigma'''$
such that
$
W(\pi'',\sigma'')=W(\pi''',\sigma''')
$
and $|\pi'''| \le C_0$
and $\Last(\pi''')=(-,-,g)$
and $\Last(\sigma''') \Last(\pi''') \Vec y$
and $\pi'''$ is domain-stable
and $\sigma'''$ is along $\pi'''$.

By definition of minimally along,
there uniquely exists the value assignment sequence $\Tilde\sigma$ such that
$\Last(\Tilde\sigma) \min(\Last(\pi''')) \Vec y$
and $\Tilde\sigma$ is minimally along $\pi'''$.

We have
\[
W(\pi''',\sigma''')\ge_\Lex W(\pi''',\Tilde\sigma)
\]
since
by definition of minimum, we have
$\Vec v_i \ge \Vec{\Tilde v_i}$ for all $i \in [1,n]$.

Define $\Hat\pi$ and $\Hat\sigma$ 
as those that give the minimum
in the body of the definition of $\rho(g,\Vec y)$.
Then 
$
\rho(g,\Vec y)=W(\Hat\pi,\Hat\sigma)
$
and
$
\Last(\Hat\pi) = (-,-,g)
$
and
$
\Last(\Hat\sigma) \min(\Last(\Hat\pi)) \Vec y$
and
$\Hat\pi$ is a domain-stable SC compositions and
$\Hat\sigma$ is a value assignment sequence minimally along $\Hat\pi$
and $|\pi| \le C_0$.

By definition of $\rho$, we have
$
W(\pi''',\Tilde\sigma) \ge_\Lex W(\Hat\pi,\Hat\sigma).
$
Hence
$\rho(f,\Vec x) >_\Lex \rho(g,\Vec y)$.
$\Box$

Note that the ranking function synthesis in this theorem
is done in intuitionistic logic.

